\section{A single proposal with bounded mean excess}
\label{sec:lookahead-cubic}

Write $W(G)=\per G$, $m_b=\mathbb E_{U_b}W$, where $U_b$ is the
uniform simple $b$-regular bipartite board, and put
$\beta_b=b(1-1/b)^{b-1}$ for $b\ge2$. For a fixed target square $S$, let
$T_b$ be the union of a uniform $b$-subset of its rows and
$\mu_b(S)=\mathbb E\log W(T_b)$. All logarithms are natural.
We use the cutoff $D=\lfloor n^{1/200}\rfloor$ from Section~2. It suffices to
describe large $n$; finitely many smaller orders may use exhaustive sampling.

Use one uniform hidden ordering of physical row labels. Begin with the
bounded polygon head, then a two-row block and a three-row block, both described below.
Next place exactly one marked block of each size $b=4,\ldots,D$, in increasing
order. The total number of previously placed rows is $q=O(D^2)$ throughout
this prefix. Continue with ordinary uniform matchings until just three rows
remain, and use the cubic sampler below. Collisions in a raw block or
infeasible masked choices are null outcomes, \emph{not} reasons to repeat the
block. Write $T$ for the ratio of reverse probability to forward subprobability
of a supported history of a given target. All reverse marks are made after
choosing the hidden row order. Different $b$-sets used in the same history
need not be independent.

\subsection{Lookahead in a marked block}

Choose an auxiliary matching $M$ first. Partition its domain sites among
the $b$ distinct physical colors, with owner counts $(v_1,\ldots,v_b)$
distributed proportionally to $\prod_j(n-v_j)!/v_j!$ on $\sum v_j=n$;
conditional on counts use the uniform assignment. For a raw block, choose
$M$ uniformly from all permutations and fill each physical color's
complement uniformly. Every legal decorated block then has forward mass
$1/Z_b$, with the integer
\begin{equation}\label{eq:v2-Zb}
 Z_b=(n!)^2\sum_{\sum_j v_j=n}\prod_{j=1}^b\frac{(n-v_j)!}{v_j!}.
\end{equation}
Indeed the owner-assignment probability is $\prod_j v_j!/n!$, and the
uniform free fills have product probability $\prod_j 1/(n-v_j)!$.
A convolution over partial sums computes the normalizer and draws the
counts; integerizing the single-color coefficients by $n!$ suffices.
In the reverse experiment, draw $M$ uniformly from the union of the
$b$ real rows. On a simple block, the owner of each edge is unique.

Take $b_0$ a sufficiently large fixed integer. For $4\le b<b_0$ use raw
blocks with a conflict veto. All these slots, including their predecessors,
are bounded in number independently of $n$. At $b\ge b_0$ draw $M$
uniformly in the current host, hence avoiding the $q$ preceding physical
rows. At color $j$, its $v_j$ owned edges are fixed; write $X_j,Y_j$
for the free domain and range, $s_j=n-v_j$, and $R_j=X_j\times Y_j$.
In this box forbid preceding physical edges and the other reserved mark
edges. If $\min_j s_j<n/2$, take uniform matching fills in these masks.
Otherwise assign the rational weight
\[
 \phi_j(F)=\prod_{e\in F}\prod_{\substack{k>j\\ e\in R_k}}(1-1/s_k)
\]
to an allowed matching $F$ of the free box, and draw proportionally.
This is a forward choice depending only on information already drawn.
The reverse experiment is unchanged. Write $P_j$ for the total weight in
the $j$th mask. On a balanced supported history, the incremental log
reciprocal \emph{relative to using ordinary matching draws in these same
slots} is
\begin{equation}\label{eq:v2-inc}
 \log\frac{Z_b}{(n!)^b}-\log W(T_b)
 +\log\frac{W(B_q)}{n!}
 +\sum_{j=1}^b\left[
     \log\frac{P_j}{s_j!}-\log\phi_j(F_j)
                         -\log\frac{W(B_{q+j-1})}{n!}\right].
\end{equation}
Here $B_i$ means the host after $i$ physical rows; the auxiliary matching
does not consume a row. Formula \eqref{eq:v2-inc} follows by replacing
the raw proposal atom for the mark and each free fill by its actual
probability. A dense unweighted slot has inverse probability $W(B_i)$.

\begin{lemma}[Weighted dense first order]\label{lem:v2-weighted-dense}
For an $s$-by-$s$ matrix $0\le w_{ij}\le1$ whose deficit $1-w$ has
maximum row or column sum $K=o(s)$,
\[
 \log(\per(w)/s!)=-\frac1s\sum_{ij}(1-w_{ij})+O(K^2/s).
\]
\end{lemma}
\begin{proof}
Interpolate $W_t=1-t(1-w)$ for $0\le t<1$. In the matching measure of
any such positive board, an edge marginal is at most $1/(s-2K)$:
from a matching containing the edge, swap its column with every other
row. The sum of relative new weights is at least $s-1-2K$ by the
row and column deficit bounds, since the removed weights are at most
one. Every new matching has a unique predecessor. This also bounds
marginals in a one-row/column minor with $s$ replaced by $s-1$.
Let $Z_{iy}$ denote the minor removing row $i$ and column $y$.
Change a matching in that minor into one missing column $z$ by sending
the row presently owning $z$ to $y$. The expected relative weight is
at least $1-K/(s-1-2K)$ by the marginal bound. Interchange $y,z$ for
the converse. Row expansion gives
$Z_{iy}/\per(W_t)=(1+O(K/s))/s$. Integrating the logarithmic derivative
$-\sum(1-w_{iy})Z_{iy}/\per(W_t)$ and taking $t\uparrow1$ proves the
claim, since the total deficit is at most $sK$.
\end{proof}

In $R_j$ the hole set is the disjoint union of exterior rows $Q$, the
$s_j$ cells of $M$ in the box, and all earlier \emph{free} fill edges
$F_i$ inside $R_j$. Earlier \emph{owned} edges are already counted in
$M$. The row and column weight deficits are at most $q+3b+1$.
Expanding the product in $\phi_j$ and applying
Lemma~\ref{lem:v2-weighted-dense} gives
\begin{align*}
 -\log(P_j/s_j!)={}&
 \frac{e_Q(R_j)}{s_j}+1+
 \sum_{i<j}\frac{|F_i\cap R_j|}{s_j}+
 \sum_{k>j}\frac{|R_j\cap R_k|}{s_js_k}
 +O((q+b+1)^2/n).
\end{align*}
Indeed the first-order soft-weight deficit is at most $2b/n$ per cell;
extending its sum to hard holes costs $O(b(q+b+1)/n)$ after division
by $s_j$. The product's quadratic remainder costs $O(b^2/n)$.
Also $-\log\phi_i(F_i)=\sum_{j>i}|F_i\cap R_j|/s_j+O(b/n)$.
Thus \emph{all} the terms recording earlier free-fill edges cancel in
\eqref{eq:v2-inc}. Conditional on the current target block and its
mark, its $q$ exterior colors are uniform among the other $n-b$ colors.
Consequently
\[
 \mathbb E[e_Q(R_j)\mid T_b,M,\text{internal order}]
      =\frac{q}{n-b}\bigl(s_j^2-e_{T_b}(R_j)\bigr).
\]
Since $e_{T_b}(R_j)\le b s_j$ and $\sum s_j=(b-1)n$, summing after
division by $s_j$ gives $q(b-1)+O(qb^2/n)$. This calculation does
not condition on the identities or decorations of earlier blocks.
The regular terms $\log[W(B_i)/n!]$ equal $-i+O(i^2/n)$ by the lemma.
The $bq$ exterior charge cancels exactly to first order.

For $x_i=v_i/n$ the remaining collision charge is
\[
 \Lambda=b+\sum_{i<j}\frac{|R_i\cap R_j|}{s_i s_j}
       = b+\sum_{i<j}\frac{(1-x_i-x_j)^2}{(1-x_i)(1-x_j)}.
\]
Here the owner sets partition each shore. Expanding the summand as
$1-x_i-x_j-x_ix_j+x_i^2x_j^2/[(1-x_i)(1-x_j)]$ proves
\begin{equation}\label{eq:v2-lambda}
 \Lambda\ge \binom b2+\tfrac12+\tfrac12\sum_i x_i^2
          \ge \binom b2+\tfrac12+\frac1{2b}.
\end{equation}
The charge in \eqref{eq:v2-inc}, apart from its first two terms, is
therefore at most $-1/2-1/(2b)+O(b(q+b+1)^2/n)$ in reverse mean. This
remains valid with a negligible additive error when including unbalanced
marks. Indeed, on \emph{any} $b$-regular block the probability of at
least $k=\lceil n/2\rceil$ owned edges of a specified color is at most
$\binom nk u_b^{n-k}/\beta_b^n$, where $u_b=(b!)^{1/b}$: prescribe
$k$ edges, use Br\'egman on the minor and Schrijver on the denominator.
As $u_b/\beta_b=\exp(O(\log(b)/b))$, $u_b\ge b/e$, this probability,
times $b$, is exponentially small uniformly for $b\ge b_0$. On a
reverse-supported unbalanced history every positive integer matching
count and every factorial has polynomial-size logarithm. One applies
the exceptional-event estimate only to the nonnormalizer terms, retaining
the exact $-\mu_b(S)$.

\subsection{One more term in the composition normalizer}

We record the precision needed to sum \eqref{eq:v2-lambda}. For every
fixed $b\ge3$ one has $Z_b/(n!)^b=O_b(\sqrt n\,\beta_b^n)$.
Indeed with $p_j=v_j/n$ the exponential rate in \eqref{eq:v2-Zb}
is $\sum_j[(1-p_j)\log(1-p_j)-p_j\log p_j]$, uniquely maximized
at the uniform point. Pairwise smoothing is strict for unequal
coordinates whose sum is less than one; faces supported on at most two
coordinates have rate zero, less than $\log\beta_b$. The constrained
Hessian at the maximum is negative definite. Stirling gives the prefactor
$O_b(n^{(2-b)/2})$ there; Gaussian summation over $b-1$ coordinates gives
$O_b(n^{(b-1)/2})$. Elsewhere the polynomial prefactor is dominated by
exponential loss. This suffices for the finitely many raw blocks.

Uniformly for $b_0\le b\le D$ the sharper estimate is
\begin{equation}\label{eq:v2-sharp-Z}
 \log[Z_b/(n!)^b]= n\log\beta_b+\tfrac12\log(2\pi n)
 +O\left(b^{-3/2}+\frac{b\ell^2}{n}+
                 \frac{\ell^3}{\sqrt{bn}}\right).
\end{equation}
To see this, let $(N_j)$ be multinomial with $n$ draws and equal cell
probabilities $1/b$, $\mu=n/b$, $U=\sum(N_j-\mu)^2/n$ and
$\lambda=b/[2(b-1)]$. The exact identity is
\[
 Z_b/(n!)^b=n!(b/n)^n\,\mathbb E\prod_j n^{N_j}/(n)_{N_j},
\]
where $(n)_t=n!/(n-t)!$. On $\max|N_j-\mu|\le\sqrt\mu\,\ell$, expand
$F(x)=x+(1-x)\log(1-x)$ about $1/b$ in the Stirling formula for this
product. Its log equals
\[
 bnF(1/b)+(b/2)\log(1-1/b)+\lambda U
             +O\bigl(b\ell^2/n+\ell^3/\sqrt{bn}\bigr).
\]
Here both linear sums vanish, $F''(1/b)=b/(b-1)$, and $F'''$ is bounded
on $[0,1/2]$. To control the complement, for every fixed $t>0$,
$\mathbb E e^{tU}=O_t(1)$ if $b_0$ is sufficiently large. Split positive
and negative deviations by Cauchy--Schwarz; negative association of
multinomial counts bounds the increasing products and the decreasing
products separately by the products of their one-dimensional expectations.
For either sign of a binomial deviation $Z=N_j-\mu$, integration of the
Chernoff bound $\Pr(|Z|\ge z)\le 2\exp[-z^2/((2+B)\mu)]$ for
$z\le B\mu$ gives $\mathbb E e^{2t(Z_\pm)^2/n}\le1+O_t(1/b)$ on
this range when $b$ is large. For $z>B\mu$, take $B$ large so the
upper tail is at most $\exp[-(2t+2)z]$; the negative tail stops at
$\mu$, and $z^2/n\le z$. The exceptional contribution is negligible
since $\mu\gg\ell^2$. For the original factorial product with all
$N_j\le n/2$, use its Riemann upper bound $\sum nF(N_j/n)\le
bnF(1/b)+U$, then Cauchy--Schwarz and the central-band tail
$O(b e^{-c\ell^2})$. If any $N_j>n/2$, use $\sum F(N_j/n)\le1$
and probability at most $b(4/b)^{n/2}$. These bounds justify the same
expansion after averaging.

For completeness $\mathbb EU=1-1/b$ and
$\Var(U)=2(b-1)(1-1/n)/b^2$: write
\[
 U-\mathbb EU=\frac2n\sum_{i<j}(\mathbf1_{\{B_i=B_j\}}-1/b)
\]
with independent uniform letters $B_i$. Different pairs are uncorrelated.
The sixth centered moment is $O(b^{-3})$. In the six-edge expansion a
draw-vertex appearing just once kills the expectation. For $v\le6$
vertices and $c$ connected components, expansion of the remaining centered
indicators bounds the absolute expectation of each product by
$C b^{-(v-c)}$: the rank of selected equality constraints plus the number
of omitted edges is at least the full rank. Multiplication by
$O(n^{v-6})$, using $v-c\ge v/2$, gives the claim. With the exponential
moment bound, Cauchy--Schwarz bounds the third Taylor remainder. Thus
\[
 \log\mathbb E e^{\lambda U}
    =\tfrac12+\frac{1-1/n}{4(b-1)}+O(b^{-3/2}).
\]
The $1/b$ term cancels the one in
$(b/2)\log(1-1/b)=-1/2-1/(4b)+O(b^{-2})$. This proves
\eqref{eq:v2-sharp-Z}, including absence of an unspecified $O(1/b)$.

The sparse one-planted-matching estimate gives, uniformly in this range,
\[
 \log m_b=n\log\beta_b+\tfrac12\log(2\pi n)-\tfrac12
                 +\tfrac12\log(1-1/b)+O(b^3/n+1/n).
\]
Combining it with \eqref{eq:v2-lambda} and \eqref{eq:v2-sharp-Z}, the
increment at grade $b$ is at most $\log m_b-\mu_b(S)+e_b$, where
$\sum_{b=4}^D e_b=O(1)$: the terms $-1/(2b)-\tfrac12\log(1-1/b)$
are $O(b^{-2})$, $\sum O(b^{-3/2})=O(1)$, and the finite-size errors
sum to $o(1)$ since $q=O(D^2)$. At the finitely many initial raw grades
their normalizer bound, constant $q$, and missing-slot comparison give
$e_b=O_{b_0}(1)$. No independence of a front block and a suffix residual
is used in cancelling the two occurrences of $\mu_b(S)$.

\subsection{The last three rows and a double cubic head}

On a cubic board $G$ the complement of a first matching $M$ consists of
$K(M)$ even cycles; it has $C(M)=2^{K(M)}$ ordered two-row completions.
Put $A=n$, $w(M)=\min(C(M),A)$ and $J(G)=\sum_M w(M)$. Select a
uniform matching and accept with probability $w(M)/A$, repeating
\emph{this matching draw} until accepted. Color the complementary cycles
by independent fair orientations. The probability of a specified ordered
factorization $z$ beginning with $M$ is exactly
\begin{equation}\label{eq:v2-cap-cubic}
 p_G(z)=\frac{1}{J(G)\max(1,C(M)/A)}.
\end{equation}
The expected number of uniform-matching calls is at most $A$. Use this
procedure for the final three rows.

For the additional \emph{pair head}, draw a first uniform permutation and
a relative derangement $\pi$ with weight
$\max(1,2^{\#\mathrm{cycles}(\pi)}/A)$. Its normalizer divided by
$(n!)^2$ is at most $1+(n+1)/A=O(1)$ since
$\sum_{\pi\in S_n}2^{\#\mathrm{cycles}(\pi)}=(n+1)!$.
The integer cycle recurrence by the cycle of the least remaining point,
stratified by the number of cycles, implements the draw with rational
probabilities. Veto all earlier conflicts. The positive overflow log in
\eqref{eq:v2-cap-cubic} for the actual last two rows cancels in reverse
mean with the negative overflow log from this head, since the two
ordered row pairs have the same marginal law.

The additional \emph{cubic head} draws two ordered 3-factorizations of
the same raw edge multiset. Each of the $n$ vertices on each shore has a
permutation $\sigma\in S_3$ identifying its three physical colors with
three auxiliary colors. Require identical left/right counts
$h_{ij}$ of occurrences of $\sigma(i)=j$. Match the two sets of
$h_{ij}$ vertices bijectively and independently for every $(i,j)$.
Veto multiple edges and collisions with previous rows. Every simple
legal pair has exactly one representation. The unrestricted count is
\begin{equation}\label{eq:v2-cubic-Z}
 Z_3^\star=\sum_h d_n(h)^2\prod_{ij}h_{ij}!,\qquad
 d_t(h)=[x^h](\per x)^t,\quad x\text{ a }3\text{-by-}3\text{ matrix}.
\end{equation}
Uniform raw generation uses the displayed weights and two independent
backward samplers from $d_t(h)=\sum_{\sigma\in S_3}d_{t-1}(h-P_\sigma)$.
Only the upper-left $2$-by-$2$ counts need be stored since margins equal
$t$. This is fixed-dimensional integer arithmetic. On a given physical
cubic block $G$, draw its reverse auxiliary coloring by
\eqref{eq:v2-cap-cubic}. Its density is at most $1/J(G)$, so the
head log reciprocal pays $-\mathbb E\log J(G)$, precisely the mean
positive cost of the terminal cubic sampler.

\begin{lemma}[The raw cubic count]\label{lem:v2-cubic-moment}
 $Z_3^\star/(n!)^3=O(n(4/3)^n)=O(m_1m_2m_3)$.
\end{lemma}
\begin{proof}
For a $3$-by-$3$ matrix $B$ of independent standard complex Gaussians,
orthogonality of monomials gives $\mathbb E|\per B|^{2n}=Z_3^\star$
by \eqref{eq:v2-cubic-Z}. Independent row-radial integration contributes
$[(n+2)!/2]^3$. For three independent unit complex rows $z,u,v$, put
$P_i=u_jv_k+u_kv_j$, $\{i,j,k\}=\{1,2,3\}$. Conditional averaging over
$z$ gives exactly $2\|P\|^{2n}/[(n+1)(n+2)]$ by the
$\operatorname{Beta}(1,2)$ squared-coordinate law. Put $t=\langle u,v\rangle$.
Expansion gives
\[
 \frac43-\|P\|^2=\frac{1-|t|^2}{3}
                 +2\sum_i|\overline{u_i}v_i-t/3|^2
 \ge c\left[1-|t|^2+\sum_i(|u_i|^2-1/3)^2\right].
\]
For the inequality rotate $v$ so $t\ge0$; the two vectors of deviations
on the two sides differ in norm by at most
$\|u-v\|+(1-t)/\sqrt3=O(\sqrt{1-t^2})$.
Conditional on $u$, $|t|^2\sim\operatorname{Beta}(1,2)$; the integral of
$\exp[-cn(1-|t|^2)]$ is $O(n^{-2})$. Squared coordinates of $u$ are
uniform on the two-dimensional simplex, so the integral of
$\exp[-cn\sum_i(|u_i|^2-1/3)^2]$ is $O(n^{-1})$. The total angular
factor, including the first $z$-integral, is $O((4/3)^n n^{-5})$.
Multiplying the radial factor proves the first assertion. For the second,
$m_1=1$, $m_2=\Theta(\sqrt n)$ by degree-two cycle enumeration, and
$m_3=\Theta(\sqrt n(4/3)^n)$ by the sparse one-plant estimate.
\end{proof}

\subsection{Mean coverage, common scale and factor cost}

\begin{proposition}[Uniform mean coverage]\label{prop:mean-cover}
For the single decorated proposal just constructed, every fixed labeled
target has $\mathbb E_{\rm rev}\log T\le\sum_{d=1}^n\log m_d+O(1)$,
with an absolute additive constant (after fixing the degree exponent
and analytic truncation).
\end{proposition}
\begin{proof}
By Lemma~\ref{lem:head}, the total high-degree contribution, including the marked polygon, is at most
 $\sum_{d>D}\log m_d+O(1)$, using the common
first-order saddle correction before summing the errors; an $O(1/d)$
unspecified remainder would not suffice. Compare the new first five
fixed head slots against missing ordinary dense slots at a cost $O(1)$:
each $W(B_i)\ge n!(1-i/n)^n$ with bounded $i$. For each $4\le b\le D$
the ordinary suffix has mean log contribution $\mu_b(S)$ and the new
head block's mean comparison to the same physical slots is at most
$\log m_b-\mu_b(S)+e_b$. Sum using the bound above. For $b=1,2,3$
use the paired cubic and overflow cancellations and
Lemma~\ref{lem:v2-cubic-moment}. These are sums of expectations under
a single normalized reverse law; no favorable intersection of events
is required.
\end{proof}

In particular a sufficiently large fixed multiple of any computable
scale $X_n\asymp\prod_d m_d$ suffices for the soft cap. The
constant-loss incidence lower gives $H/|\mathcal L_n|=O(1)$ for such
a choice. This does not assert a constant pointwise reverse gate
probability. There is a counting consequence independent of the cost
of the gate: the forward marginal on $S$ is
$Q(S)=\mathbb E_{\rm rev}(1/T)\ge\exp(-\mathbb E_{\rm rev}\log T)$.
Since $\sum_S Q(S)\le1$ and the proposal is the same for every $S$,
one obtains $|\mathcal L_n|=O(\prod_d m_d)$. Together the bounds give
$|\mathcal L_n|=\Theta(X_n)$, not an asymptotic equivalent.

Finally, the weighted fill does not require a larger matching graph.
For each future $k$ let $l_k=\max(0,|X_j\cap X_k|+|Y_j\cap Y_k|-s_j)$.
If the chosen matching has $m_k'$ cells in $R_k$, then
$0\le m_k'-l_k\le v_k$. Thus
$\psi(F)=\phi_j(F)/\prod_{k>j}(1-1/s_k)^{l_k}$ lies between $e^{-4}$
and $1$ on balanced boxes: use $\sum v_k=n$ and $s_k\ge n/2$.
Rejecting a uniform matching with this acceptance rule gives the
weighted draw in constant expected attempts. Its normalizer is an
ordinary unweighted matching count times a known rational times the
mean of a $\psi$-coin of bias at least $e^{-4}$. There is one such
scalar per physical fill, $O(D^2)$ in total. For each of the two cubic
boards the only new scalar is the bias
$\theta=J(G)/(A\,W(G))=\mathbb E_M[w(M)/A]\ge1/n$.
These positive factors and their reciprocals use the same
geometric/negative-binomial gate as the matching marginals, without
anchor generation for the scalar itself. No unknown $J$ or weighted
permanent is evaluated to decide a coin.
